% 引言
\chapter{引言}\label{chap:intro}

\section{研究背景与问题}

硅interposer和先进封装技术的持续进步，推动晶圆级集成从学术概念进入工业量产。
TSMC的InFO-SoW平台在Tesla Dojo训练Tile中集成了25个D1裸片~\cite{tesla2022dojo}；
Cerebras WSE-3以单晶圆级芯片集成约90万个处理核心，通过2D Mesh实现20 PB/s的片上聚合带宽~\cite{cerebras2022wse2}。
在此趋势下，一个直接的架构问题是：能否在晶圆上构建一个网络交换机？

晶圆级平台在该方向上具备两项天然优势。
第一，硅interposer提供的裸片间互联密度远超传统PCB——
以$12 \times 12$\,mm裸片、45\,$\mu$m pitch的$\mu$bump阵列计算，单裸片可布置约7万个互联点，
聚合IO带宽可达Pbps量级~\cite{ucie2.0-2024}，为超大radix交换网络提供了物理基础。
第二，interposer上的互联协议灵活可配——
UCIe Advanced（2.5D）每bit功耗低至0.25\,pJ，SerDes VSR/MR/LR覆盖10\,cm至100\,cm的reach范围，
光学互联（如Ayar Labs TeraPHY）则可提供8\,Tbps的单chiplet带宽和$<$5\,pJ/bit的能效~\cite{oif-cei-5.1}——
设计者可根据裸片间距自适应选择速率和功耗等级。

学术界对此方向已做出响应。
Chen等人~\cite{chen2024waferscale}首次量化了晶圆级交换机的radix提升潜力，
指出仅面积约束下可达传统芯片级交换机32倍的端口密度；
Feng和Ma~\cite{feng2024switchless}提出”无交换机Dragonfly”架构，
利用晶圆级高密度互联替代传统高radix物理交换机。
Wan等人~\cite{wan2024architectural}在2D Mesh物理拓扑上叠加Butterfly Fat-Tree逻辑拓扑，
实现了896端口、8.96\,Tb/s的晶圆级交换系统；
Yang等人~\cite{yang2025pd}提出的TickTock框架则面向LLM训练场景
进行物理-逻辑拓扑的联合协同设计。
这些工作确立了晶圆级交换机作为一个独立研究方向的价值，
然而它们也共同揭示了一个尚未被充分解决的核心困难：
晶圆级交换机的设计空间极大，且体系结构决策与多类物理约束之间存在深度耦合。

\section{晶圆级交换机设计的两个根本困难}

\subsection{设计空间组合爆炸}

晶圆级交换机需在拓扑族（Dragonfly、Mesh、Torus、k-ary n-cube等）与拓扑参数（Dragonfly的$a,p,h,g$）~\cite{kim2008technology}、
路由策略（Valiant、minimal、自适应）、互联标准（UCIe Advanced/Standard、SerDes VSR/MR/LR、光学互联）~\cite{ucie2.0-2024,oif-cei-5.1}、
裸片规格（面积、功耗、bump pitch）、封装方式和冷却方案等维度上做出联合决策。
这些选择跨越体系结构到物理实现多个抽象层次，笛卡尔积超过$10^6$个候选设计点。
在此规模下，cycle-accurate仿真器单次评估需数分钟量级~\cite{dally2004principles}，完全不可行；
chiplet设计空间探索工具（如FPIA~\cite{jiao2024fpia}、RapidChiplet~\cite{iff2025rapidchiplet}、
FireLink~\cite{li2025firelink}）虽能提供快速评估，但各自聚焦于单一物理维度（布局布线或成本建模），
缺乏跨层次的联合决策能力。

针对这一困难，本文提出两层设计空间探索架构：
外层负责离散技术选型的枚举与剪枝，内层通过线性规划对候选设计点进行快速可行性判断，
将单次评估从分钟级压缩至秒级。

\subsection{体系结构与物理约束的深度耦合}

设计空间巨大是困难的表象，更深层的问题在于体系结构与物理约束之间的闭环耦合。
在传统体系结构设计中，性能分析与物理实现可以大致分离——
先确定微架构和拓扑，再交由物理设计团队处理布局、供电和散热。
然而在晶圆级交换机中，这一分离不再成立。

拓扑和路由决策决定了每条链路上的流量分布，
流量分布决定了每条链路需要的物理lane数量和信号$\mu$bump占用~\cite{ucie2.0-2024}，
lane的物理层功耗决定了裸片总功耗，
功耗的空间分布通过热传导决定了温度梯度和翘曲风险~\cite{mfit2025}，
而电源bump需求反过来挤占信号bump的可用面积~\cite{electrical-thermal-chiplet2024}——
四个环节构成闭环。任何单一维度的优化都可能被另一维度的约束推翻，
不存在“先检查A再检查B”的序贯验证流程。

现有工具各自处理这一耦合链中的单个环节——
性能仿真器处理流量分布，bump预算工具处理信号完整性，热分析工具处理温度和翘曲——
但没有任何工具能够在整个耦合链上联立判断综合可行性。
本文将耦合链上全部环节统一表示为关于同一组变量的线性约束，
构建联立求解框架，将多轮独立验证替换为单次综合判定。

\section{核心洞察：两类约束与两层架构}

上述困难指向同一个缺口：
晶圆级交换机缺乏一个早期设计决策工具，能够同时考虑全部物理约束、
快速判定候选方案的可行性、并给出明确的设计改进方向。

构建这一工具的前提是认识约束的结构特征。
对晶圆级交换机涉及的物理、几何和工程约束进行系统分析后，
我们发现：尽管约束数量庞大且物理来源各异，它们可明确归为两类。

\textbf{第一类为离散选择}，包括互联标准（UCIe Advanced或Standard~\cite{ucie2.0-2024}）、
裸片布局、工艺节点等，决定了系统的拓扑结构和物理参数。
这类选择的数学性质是枚举型的，求解依赖于启发式搜索或智能选型策略。

\textbf{第二类为带宽需求驱动的近似线性约束}，
包括bump占用（链路负载 $\times$ 每bit所需lane数）、
功耗（互联标准的功耗估计 $\times$ 链路负载）、
温度梯度（热传导方程中热源项的线性叠加）和翘曲温差。
这类约束的数学性质是线性不等式——对带宽需求的依赖是线性的或可通过保守估计线性化。

这一分类将设计空间探索问题分割为两个不同层面的子问题。
上层为离散枚举与智能选型，通过启发式搜索或现有chiplet设计空间探索工具产生候选设计点。
下层为多约束耦合线性规划，将全部线性约束表示为不等式组联立求解，
判定候选设计在全部物理约束下的综合可行性。
两层通过物理参数接口解耦：上层参数确定后，下层所有系数矩阵随之固定。

基于这一架构，本文提出WaferDSE——面向晶圆级交换机早期设计空间探索的统一线性规划框架。
框架将体系结构决策与多物理约束的耦合判断，
从“依赖经验的独立验证”提升为“基于数学规划的联立求解”。

\section{本文贡献}

\begin{enumerate}
    \item \textbf{两层设计空间探索架构。}
    针对晶圆级交换机设计空间超过$10^6$且体系结构与物理约束深度耦合的困难，
    提出外层（离散枚举与智能选型）与内层（多约束线性规划）结合的两层架构。
    内层将性能、几何、功耗三族物理约束统一为同一变量上的线性不等式组，
    每组约束支持多种建模精度，精度升级仅替换系数矩阵而不改变框架结构，
    将最困难的跨维度耦合判断封装为单个可高效求解的线性规划。

    \item \textbf{基于灵敏度分析的设计指导框架。}
    现有设计空间探索方法仅回答“可行或不可行”，架构师无法获知哪个约束是瓶颈、
    改进方向在哪里、改进收益有多大。
    本文利用线性规划对偶变量提取每个物理约束的影子价格，
    并推导物理带宽天花板的解析公式，
    使设计迭代获得明确的优先级排序依据和定量的改进收益评估。

    \item \textbf{设计空间探索实验与结构发现。}
    在Dragonfly拓扑家族上运行大规模设计空间探索扫描，
    通过批量线性规划求解与$B_{\max}$分析的交叉验证评估候选设计的质量。
    实验识别出若干具有实际潜力的设计结构，
    同时揭示了晶圆级交换机物理瓶颈的定量缩放规律，
    为拓扑扩展和技术路线选择提供数据支撑。
\end{enumerate}
